\section{Related Work}
\label{sec:related}

\subsection{Robotic Inspection of Electrical Substations}
Autonomous ground robots are increasingly deployed to inspect electrical
substations, reducing human exposure to high-voltage hazards while improving
the regularity of meter reading, switch-status recognition and infrared
thermography~\cite{li2024state}. Navigation in this domain must cope
with dense equipment layouts and safety-critical clearances; recent work
improves path planning, e.g., by combining deep reinforcement learning with
ant-colony optimization for substation routing~\cite{zhang2025substation},
reporting gains over deep-RL and classical baselines such as PPO, DDQN, and
A$^{\!*}$; broader surveys catalogue perception and path-planning techniques for
substation mobile robots~\cite{substationpp2024review}. Contemporary substation
systems increasingly couple deep learning with the navigation stack---for example
attention-enhanced semantic segmentation to guide the inspection
robot~\cite{substation2025segnav}, and deep-learning inspection pipelines for
ultra-high-voltage sites~\cite{substation2024uhvdl}. These systems, however,
treat the round as a fixed coverage problem: the robot follows planned routes and
the perception pipeline classifies what is seen.
\emph{In contrast}, we add a reactive affective response \emph{during} the
round, letting saliency steer the robot toward anomalies and away from hazards.

\subsection{Reactive Navigation: Potential Fields and DWA}
Classical reactive methods react to local sensing without a learned value
function. Artificial potential fields (APF) sum attractive and repulsive
forces~\cite{khatib1986} but are well known to stall in local
minima~\cite{koren1991}; the Dynamic Window Approach selects feasible
velocity commands by short-horizon prediction~\cite{fox1997}. Our affective
field $\Phi$ is reminiscent of attraction/repulsion, but the reactive response is
issued by a threshold-triggered controller that preempts the patrol with a goal,
and the deliberative layer prevents the global stalls that afflict pure
potential-field control. We make this comparison concrete by evaluating a
conventional artificial-potential-field controller as a baseline on the same
events (Section~\ref{sec:results}).

\subsection{Reinforcement Learning for Navigation and Reward Shaping}
Reinforcement learning underpins many navigation policies~\cite{suttonbarto2018};
deep reinforcement learning in particular now drives a broad line of map-less and
dynamic-environment navigation, surveyed with its persistent sim-to-real and
safety challenges~\cite{drlrobotics2024survey}. Within this line, \emph{reward
shaping} accelerates learning; potential-based shaping is the canonical form that
preserves the optimal policy~\cite{ng1999}, and shaped rewards have been used for
obstacle avoidance and goal reaching~\cite{zhang2020reward}. A related hybrid line lets RL \emph{tune the
parameters} of an artificial potential field online~\cite{zhang2023apf}. As we
derive in Appendix~\ref{app:update}, writing our field as an additive term in a
value update is, for a fixed learning rate, \emph{equivalent} to augmenting the
reward with a saliency signal; we therefore do not claim it as a new learning
operator, and the architecture evaluated here contains no learning component. What
distinguishes this work is not the term in isolation but the \emph{architecture}:
a saliency-driven attract/repel impulse produced by a reactive layer on top of an
active-sensing patrol, with the value-update formulation offered as a substrate
for future learned policies. The repulsive component and the deliberative
avoidance it triggers also connect to \emph{safe} reinforcement learning, whose
central concern is maintaining safety during learning-driven
control~\cite{brunke2022safe}.

\subsection{Deliberative Stacks: Nav2 and MPPI}
We build the deliberative layer on a mature, reproducible
stack~\cite{macenski2020nav2}, using a Model Predictive Path Integral
controller for path following~\cite{williams2017mppi}. Rather than proposing a
new planner or controller, we keep this layer standard and add the reactive
affective layer on top, so the contribution is orthogonal to (and benefits
from) ongoing improvements in industrial navigation stacks.

\subsection{Active Perception and Affective Robotics}
Active perception holds that perception is purposive: an agent should act to
acquire informative observations~\cite{bajcsy1988}; its modern instantiation is
\emph{informative path planning}, where the robot plans viewpoints to maximize
information gain during mapping and inspection~\cite{ipp2024nonconvex}. Emotion-
and affect-inspired robotics, in turn, models internal drives that bias
behavior~\cite{breazeal2003}. Architecturally, our two layers follow the
classical hybrid deliberative--reactive paradigm, in which a reactive behavior is
arbitrated on top of a deliberative planner~\cite{hybrid2015delibreact}; here the
arbitration is a simple priority (safety before inspection) together with a
saliency threshold. We
connect these threads operationally: the patrol is an \emph{active-sensing} layer
that brings the sensors to where events may occur, and \textsc{Stay Alert} is the
\emph{affective} layer that, once an event is within range, produces an approach
or retreat impulse. Table~\ref{tab:comparison} positions our method against
representative reactive, learning-based, and substation-specific approaches.
Within this scope, we are not aware of a prior method that couples an affective
attraction/repulsion term---formulated in the Bellman update---with an
active-sensing patrol for substation inspection.

\begin{table*}[!t]
\caption{Positioning of the proposed method against representative approaches.
``Active sensing'' denotes moving the platform to acquire observations beyond a
fixed route.}
\label{tab:comparison}
\centering
\footnotesize
\setlength{\tabcolsep}{4pt}
\begin{tabular}{p{2.3cm}p{2.4cm}p{2.7cm}p{1.4cm}p{2.5cm}p{1.7cm}p{2.2cm}}
\hline
\textbf{Method} & \textbf{Learning strategy} & \textbf{Reactive mechanism} &
\textbf{Active sensing} & \textbf{Safety handling} & \textbf{Validation} &
\textbf{Reported metrics} \\
\hline
Classical APF / DWA~\cite{khatib1986,fox1997} & None & Force summation / velocity
search & No & Repulsive force & Sim / analytical & Reach, path length \\
RL-tuned APF~\cite{zhang2023apf} & RL tunes APF gains & Potential field & No &
Repulsive force & Simulation & Success, collisions \\
Substation DRL+ACO~\cite{zhang2025substation} & Deep RL + ant colony & None
(global planner) & No & Planned clearance & Sim + field & Path length, completion
\\
Informative path planning~\cite{ipp2024nonconvex} & Planning (partly learned) &
None & Yes (viewpoints) & Task-dependent & Sim + real & Info.\ gain, coverage \\
\textbf{This work (\textsc{Stay Alert})} & Value augmentation,
threshold-triggered (not learned) & Affective FSM (attract / repel) & Partial:
fixed-patrol sensor exposure + event-triggered repositioning (no informative view
planning) & Reactive retreat + deliberative keepout & Simulation (Gazebo) &
Inspection, safety, mission time (effect sizes, CIs) \\
\hline
\end{tabular}
\end{table*}
